\documentclass[10pt,conference,compsocconf]{IEEEtran}

\usepackage{hyperref}
\usepackage{graphicx}
\usepackage{amsmath}
\usepackage{amssymb}
\usepackage{booktabs}
\usepackage{xcolor}
\usepackage{microtype}

\begin{document}

\title{Metric-Triggered WSD Decay for Language Model Pretraining}

\author{
  Pierre Bernadet, Georges-Alex Nahas\\
  \textit{CS-439 Optimization for Machine Learning, EPFL, Spring 2026}
}

\maketitle

\begin{abstract}
  Warmup-Stable-Decay (WSD) learning rate schedules have emerged as a flexible
  alternative to cosine annealing for language model pretraining, decoupling the
  stable training phase from a short, sharp decay. Yet the key hyperparameters
  governing this decay—its timing, duration, and final learning rate ratio—remain
  largely empirical and are unlikely to generalize across models and datasets.
  In this work, we study the relationship between online training signals
  (loss variance and gradient signal-to-noise ratio) and the optimal decay
  trigger in WSD, using a LLaMA-style model trained on FineWeb.
  We reproduce scheduler comparison results from prior work, analyze how decay
  timing and magnitude affect final validation loss, and propose adaptive
  metric-based policies to automatically identify the ``sweet spot'' for
  initiating decay—the point at which the stable phase yields diminishing returns.
  % TODO: fill in concrete results once figures are finalized.
\end{abstract}


% ============================================================
\section{Introduction}
% ============================================================

Cosine learning rate annealing~\cite{kaplan2020scaling} has long been the
default schedule for pretraining large language models, but it imposes a hard
token budget: the entire training duration must be fixed before training begins,
and extending a run requires restarting from scratch. WSD
schedules~\cite{minicpm} address this limitation by separating training into
three phases—warmup, a long stable phase at peak learning rate, and a short
decay—making continued training and branching straightforward.

Despite this appeal, WSD introduces its own set of design choices. When should
decay begin? How long should it last, and how deep should the learning rate
drop? Existing work~\cite{minicpm,wen2024river} provides recommendations that
are largely empirical and tuned for specific model and data scales. Finding the
right timing—the \emph{sweet spot}—is non-trivial: decaying too early
forfeits further stable-phase progress, while decaying too late may leave
potential gains on the table if the budget is fixed.

We define the sweet spot as the step at which the stable phase has converged
to a plateau: further training at peak learning rate yields diminishing
validation-loss improvements, and initiating decay will most efficiently
compress remaining progress. Our work addresses three questions:
\begin{enumerate}
  \item Can we reproduce the WSD vs.\ cosine comparison from prior
        work~\cite{wen2024river} on our model and dataset?
  \item Do online training metrics correlate with validation loss under
        different decay timings and magnitudes, allowing us to predict the
        sweet spot?
  \item Can we construct meaningful adaptive policies, based on these
        metrics, that match or improve upon fixed-schedule baselines?
\end{enumerate}


% ============================================================
\section{Related Work}
% ============================================================

Learning rate scheduling is a critical and sensitive component of LLM
pretraining~\cite{kaplan2020scaling,hoffmann2022chinchilla}.
Cosine annealing has been the dominant choice~\cite{dubey2024llama3}, but it
requires committing to a total token budget upfront, limiting training
flexibility~\cite{defazio2024road,muennighoff2023scaling}.

To overcome this rigidity, \citet{minicpm} introduced the WSD scheduler as
part of the MiniCPM training framework. WSD keeps the learning rate at its
peak value throughout a long stable phase and compresses the final descent
into a short, steep decay, enabling seamless extension of training runs.
\citet{wen2024river} provided a theoretical grounding for WSD through the
\emph{river valley} loss landscape perspective: the stable phase navigates
the model along a flat valley bottom, while decay rapidly descends to a sharp
minimum. Their analysis, conducted under SGD dynamics, also introduced
WSD-S—a smoothed variant of the decay curve—and empirically compared WSD-S
against cosine annealing, demonstrating faster convergence per token.
Further empirical evidence for WSD's efficiency has been reported in a number
of subsequent pretraining studies.

All of this prior work treats the timing and magnitude of the decay as fixed,
hand-tuned hyperparameters. Our study fills this gap by asking whether
\emph{online} training signals can replace these manual choices.


% ============================================================
\section{Materials and Methods}
% ============================================================

\subsection{Datasets}

We use FineWeb~\cite{penedo2024fineweb} (\texttt{sample-100BT}, 100B tokens)
as our primary pretraining corpus, a large-scale, high-quality web crawl widely
adopted in recent open LLM research. To test cross-dataset generalization of our
policies, we also run selected experiments on
C4~\cite{raffel2020c4}, another standard English pretraining benchmark.
Both datasets offer sufficient scale for our computation budget while remaining
representative of realistic pretraining conditions.

\subsection{Model}

We train a LLaMA-style~\cite{dubey2024llama3} causal language model with 8
transformer layers, hidden dimension 512, intermediate size 2048, and 8
attention heads, yielding approximately 40M parameters. This size provides an
acceptable trade-off between LLM representativeness and the constraint of
running many ablation experiments.

We use the GPT-2 tokenizer (vocabulary size 50,257) and a context length of
1024 tokens. Optimization uses AdamW~\cite{loshchilov2019adamw} ($\beta_1 =
0.9$, $\beta_2 = 0.95$, weight decay 0.1) with a peak learning rate of $6
\times 10^{-4}$ and gradient clipping at norm 1.0. Following~\citet{wen2024river},
we adopt AdamW despite the WSD theory being developed under SGD, as AdamW
is well established for LLM pretraining. Effective batch size is 524,288 tokens
(batch 8, accumulation 64), giving a total token budget of 5.24B tokens over
10,000 steps. We use 5,000 validation samples sampled from a held-out stream of
FineWeb, evaluated every 50 steps. Warmup lasts 100 steps.

\subsection{Experimental Setup}

All experiments share the base configuration above unless stated otherwise.
We design four experiments to progressively build from reproduction to adaptive policies.

\subsubsection{Exp.\,1 — Scheduler Comparison}
We reproduce the cosine vs.\ WSD-S comparison of~\citet{wen2024river} on our
model. A single WSD-S run (decay starting at 90\% of steps, lasting 10\%,
final LR ratio 0.1) is compared against a cosine run with identical peak LR
and training budget. This follows the paper's recommended hyperparameters
and confirms that our setup reaches the plateau phase before decay.

\subsubsection{Exp.\,2 — AdamW vs.\ Quasi-SGD Decay}
The theoretical analysis of WSD in~\citet{wen2024river} assumes SGD.
To better align with this theory during the decay phase, we shrink the AdamW
momentum coefficients to $(\beta_1, \beta_2) \approx (0, 0)$ at the moment
decay begins, effectively disabling first- and second-moment accumulation and
recovering a quasi-SGD update rule (see Appendix~\ref{app:adamw-sgd} for
the derivation). We compare three conditions: (a)~WSD with standard AdamW,
(b)~cosine with standard AdamW, and (c)~WSD with AdamW betas shrunk at
decay onset.

\subsubsection{Exp.\,3 — Decay Timing and Amount Sweep}
To map how the sweet spot depends on when and how strongly the decay is
applied, we run a grid of WSD-S runs varying the decay start step (as a
fraction of total training: 50\%, 70\%, 80\%, 90\%) and the final LR ratio
($r \in \{0.5, 0.3, 0.1, 0.05, 0.01\}$). For each configuration we record
the final validation loss and the values of two online metrics at decay onset:
\emph{loss variance} (variance of training loss over a 50-step window) and
\emph{gradient SNR} (ratio of the mean gradient norm to its standard deviation
across layers). The goal is to find reliable correlations between these
pre-decay signals and post-decay performance.

\subsubsection{Exp.\,4 — Adaptive Decay Policies}
Based on the correlations found in Exp.\,3, we propose three adaptive policies
that monitor training in real time and trigger decay automatically once the
signal indicates the sweet spot has been reached:
\begin{itemize}
  \item \textbf{Loss-variance policy}: triggers when the current
        \texttt{loss\_variance} falls below the 3rd percentile of values
        observed in the preceding lookback window (10\% of total steps).
  \item \textbf{Grad-SNR policy}: analogously triggers on the 3rd percentile
        of \texttt{grad\_snr}.
  \item \textbf{Combined policy}: requires \emph{both} conditions
        simultaneously.
\end{itemize}
All policies are gated: they cannot fire before 70\% of total steps (7,000),
to avoid premature decay. Decay length is fixed at 1,000 steps with final LR
ratio 0.1. Percentile thresholds are computed from the lookback window at
gate time (step 7,000). We compare each policy against three fixed baselines
that start decay at 70\%, 80\%, and 90\% of steps respectively—the last
being the setting recommended by~\citet{wen2024river}.


% ============================================================
\section{Results}
% ============================================================

Figures~\ref{fig:exp1}--\ref{fig:exp4} summarize the results of each
experiment in the order described above.

\begin{figure}[tbp]
  \centering
  % \includegraphics[width=\columnwidth]{figures/exp1_scheduler_comparison}
  \fbox{\parbox{0.9\columnwidth}{\centering\vspace{1.5cm}
      \textbf{[Fig.\,1: Validation loss curves for cosine vs.\ WSD-S,}\\
      \textbf{Exp.\,1 reproduction of \citet{wen2024river}]}
      \vspace{1.5cm}}}
  \caption{Cosine vs.\ WSD-S validation loss on FineWeb (10k steps).
    Decay starts at step 9,000 (90\%) for WSD-S.}
  \label{fig:exp1}
\end{figure}

\begin{figure}[tbp]
  \centering
  % \includegraphics[width=\columnwidth]{figures/exp2_quasi_sgd}
  \fbox{\parbox{0.9\columnwidth}{\centering\vspace{1.5cm}
      \textbf{[Fig.\,2: WSD standard AdamW vs.\ WSD quasi-SGD betas}\\
      \textbf{vs.\ cosine, validation loss, Exp.\,2]}
      \vspace{1.5cm}}}
  \caption{Effect of shrinking AdamW betas at decay onset to approximate
    SGD dynamics as assumed by the WSD theoretical analysis.}
  \label{fig:exp2}
\end{figure}

\begin{figure}[tbp]
  \centering
  % \includegraphics[width=\columnwidth]{figures/exp3_sweep}
  \fbox{\parbox{0.9\columnwidth}{\centering\vspace{1.5cm}
      \textbf{[Fig.\,3: Final validation loss vs.\ decay start/amount,}\\
      \textbf{and metric correlations, Exp.\,3]}
      \vspace{1.5cm}}}
  \caption{Heatmap of final validation loss across decay timing $\times$
    final LR ratio, with correlation of pre-decay loss variance and grad
    SNR to final loss overlaid.}
  \label{fig:exp3}
\end{figure}

\begin{figure}[tbp]
  \centering
  % \includegraphics[width=\columnwidth]{figures/exp4_policies}
  \fbox{\parbox{0.9\columnwidth}{\centering\vspace{1.5cm}
      \textbf{[Fig.\,4: Adaptive policies vs.\ fixed baselines,}\\
      \textbf{validation loss, Exp.\,4]}
      \vspace{1.5cm}}}
  \caption{Validation loss for adaptive policies (loss variance, grad SNR,
    combined) vs.\ fixed baselines at 70\%, 80\%, and 90\% trigger.}
  \label{fig:exp4}
\end{figure}

% TODO: fill in once figures are generated.


% ============================================================
\section{Discussion}
% ============================================================

% TODO: discuss once results are finalized.
% Tentative structure:
% - Exp1: does our model reproduce the WSD vs cosine gap?
% - Exp2: does quasi-SGD decay improve on standard AdamW WSD?
% - Exp3: which metrics correlate best with post-decay loss?
%         Does earlier decay (more budget) always lose?
% - Exp4: do policies match the best fixed baseline?
%         When do they fire vs fixed triggers?
%         Limitations: threshold sensitivity, gate choice.


% ============================================================
\section{Summary}
% ============================================================

% TODO: fill in once results are finalized.
% Tentative:
% - WSD outperforms cosine with similar settings on our model/data.
% - [result of quasi-SGD experiment].
% - Online metrics (loss variance, grad SNR) [do/do not] reliably
%   predict the sweet spot for decay timing.
% - Adaptive policies achieve [competitive/improved] results vs.
%   fixed 90% baseline, without requiring manual timing selection.


% ============================================================
\newpage
\bibliographystyle{IEEEtran}
\bibliography{references}


% ============================================================
\appendix

\section{AdamW to Quasi-SGD Beta Shrinking}
\label{app:adamw-sgd}

The AdamW update at step $t$ is:
\begin{align}
  m_t &= \beta_1 m_{t-1} + (1 - \beta_1)\,g_t, \\
  v_t &= \beta_2 v_{t-1} + (1 - \beta_2)\,g_t^2, \\
  \hat{m}_t &= m_t / (1 - \beta_1^t), \quad
  \hat{v}_t = v_t / (1 - \beta_2^t), \\
  \theta_t &= \theta_{t-1} - \eta\,\hat{m}_t / (\sqrt{\hat{v}_t} + \varepsilon)
             - \eta\,\lambda\,\theta_{t-1},
\end{align}
where $g_t$ is the gradient, $\eta$ the learning rate, and $\lambda$ the weight
decay coefficient. When $\beta_1 \to 0$, the first moment collapses to the
instantaneous gradient: $m_t \approx g_t$, eliminating gradient averaging.
When additionally $\beta_2 \to 0$, the second moment tracks
$v_t \approx g_t^2$, so the adaptive scaling $\hat{m}_t/\sqrt{\hat{v}_t}$
approaches $\text{sign}(g_t)$—a per-element signed update reminiscent of
SignSGD.

In practice, setting $(\beta_1, \beta_2) \leftarrow (0,\,0)$ at the onset of
the decay phase effectively disables the exponential moving averages built up
during the stable phase and replaces them with the current-step gradient.
This brings the optimizer closer to the pure SGD assumed in the theoretical
analysis of~\citet{wen2024river}, at the cost of discarding useful second-order
information. The empirical comparison in Exp.\,2 tests whether this theoretical
alignment translates into a measurable gain.

\section{Correlation Analysis Details}
\label{app:correlations}

% TODO: Pearson / Spearman correlation tables between pre-decay metrics
% (loss_variance, grad_snr) and final validation loss across Exp.3 grid.
% Include scatter plots and regression lines.

\section{Additional Plots}
\label{app:extra-plots}

% TODO: Per-run training curves, trigger step distributions for adaptive
% policies, C4 cross-dataset generalization results.

\end{document}
